% ID: SEL-D0035-OB-E01
% ESTADO: cotejada. Segunda pasada visual y borrador compilado revisados.
% Fuente: Selectividad/matematicas_aplicadas/titular_septiembre_EXAMEN 1-ANDALUCÍA.pdf
% SHA-256: 7a09a02641ea8de345b94f49b27515e58436e032c7401867f86dc523490a3ea5
% Página: 2; opción B; ejercicio original 1.

\begin{enumerate}[label=\alph*)]
\item \textbf{(0.8 puntos)} Represente el recinto definido por las siguientes inecuaciones:
\[
x+y\leq3\qquad2x+y\geq4\qquad y\geq-1
\]
\item \textbf{(0.25 puntos)} Razone si el punto $(2,1)$ pertenece al recinto anterior.
\item \textbf{(1.2 puntos)} Obtenga los vértices del recinto y los valores mínimo y máximo de la función $F(x,y)=5x+4y$ en ese recinto, indicando en qué puntos se alcanzan.
\item \textbf{(0.25 puntos)} Razone si la función $F$ puede alcanzar el valor 9 en el recinto anterior.
\end{enumerate}
